\chapter{The HISTORY of AREA 51}\label{ch:historyarea51}
\markboth{The History of Area 51}{The History of Area 51}

\chapquote{``Three may keep a secret, if two of them are dead.''}{Benjamin Franklin}

% ---------------------------------------------------------------
\section{Security and Freedom}\label{sec:securityfeedom}

The most remarkable thing about Area 51 is not that it kept secrets. It is that until
2013 the United States government would not concede that the place existed, while
simultaneously prosecuting people for photographing it. You cannot trespass on nothing.

Chapter~\ref{ch:roswell} left us with the question and no way to answer it from inside a single
case: whether a classification stamp tracks anything real, or merely marks a purpose the public is
not permitted to audit. Groom Lake is where that question can finally be tested, because here the
secrecy is not a wartime filing decision buried in a hangar inventory. It is written into land law,
environmental law and the rules of evidence, and all of it is on the record.

\figwrap[r]{0.38}{Commercial satellite imagery of the Groom Lake facility: runways, hangars, and the dry lakebed.}{https://commons.wikimedia.org/wiki/File:Area51_ikonos.jpg}{fig:groomlake}
The facility sits on the southern shore of Groom Lake, Nevada, inside the vast Nevada Test and Training Range (see Figure~\ref{fig:groomlake}). Its formal designation has changed repeatedly --- Watertown,
Paradise Ranch, Homey Airport, Detachment 3 of the Air Force Flight Test Center --- and
the name ``Area 51'' appears to derive from Atomic Energy Commission grid numbering. The
existence of the base was officially acknowledged in June 2013 with the release of a
redacted CIA history of the U-2 programme, and the acknowledgement was almost an
afterthought: the document simply used the name in passing, as though it had never been
controversial.

Here is the part that should interest a sceptic more than any rumour of alien hardware.
The legal machinery built to protect this base is genuinely extraordinary and entirely
documented. The airspace above it, R-4808N, is closed absolutely --- known to pilots as
``the Box''. The perimeter is patrolled by contractors, universally called ``camo
dudes'', who are authorised to use deadly force and who are not federal officers. In 1984
the government seized an additional 89,\!000 acres of public land, including Whitesides
and Freedom Ridge, specifically to eliminate the last public vantage points --- without
the congressional hearing normally required. In the 1990s, when former workers sued over
exposure to burning toxic waste, the government invoked the state secrets privilege to
have the case dismissed, and President Clinton signed an annual exemption from
environmental disclosure law that has been renewed by every president since.

\begin{funfact}
The workforce commutes daily by air. A fleet of unmarked Boeing 737s operates out of a
private terminal at Las Vegas Harry Reid International under the callsign ``Janet''. You
can watch them take off from the airport's public parking garage, and the airline is dealt
with in full in Section~\ref{sec:janet}. A secret you can photograph on a schedule is a
strange sort of secret.
\end{funfact}

I want to be precise about what this proves and what it does not. It proves the state
will lie, at scale, for decades, about the mere existence of a place, and will restructure
land law and environmental law to do it. That is established fact, not conjecture. It does
not prove what the lie was concealing. The documented answer --- reconnaissance and
stealth aircraft --- is entirely sufficient to explain the secrecy, and Section~\ref{sec:burdenproof}'s
razor applies with full force here.

But notice the cost. Every genuine denial the government issued about Groom Lake was
false, and each false denial purchased a little more public willingness to believe the
next extraordinary claim. Institutional secrecy manufactures conspiracy theories as a
by-product. You cannot spend fifty years teaching a population that official denials mean
nothing and then act surprised when they stop believing the true ones.

Which brings the Roswell question to its answer, and the answer is worse than a straight yes or no.
Top Secret does mean something --- Groom Lake really did hold the U-2, the Blackbird and the
stealth programmes, and every one of them was worth hiding from the Soviet Union. The stamp was
honest about the aircraft. It was simultaneously used to seize public land without a hearing, to
have poisoned workers thrown out of court, and to deny the existence of a facility visible from a
public hill. The same instrument did both jobs, and no citizen outside the fence could tell which
job it was doing on any given day. That is the real cost of secrecy, and it is not the loss of any
particular secret. It is the loss of the public's ability to distinguish a legitimate one from a
convenient one --- and once that is gone, freedom and security stop trading against each other and
simply decline together.

% ---------------------------------------------------------------
\section{The History of Aviation}\label{sec:aviationsecrecy}

Section~\ref{sec:historyaviation} gives you the timeline of flight --- balloons to Blackbirds, and how brutally fast it
happened. I am not going to repeat it here. What this section is about is a different history that
runs alongside it and explains why a place like Groom Lake exists at all: the history of \emph{where}
aviation got developed, and how flight testing went from a public spectacle to a state secret in
about forty years.

Because it started completely open. Early aviation was a performance industry. The Wrights
demonstrated at Fort Myer and at Le Mans in front of crowds and newspapermen. Between 1909 and 1914
the fastest aircraft in the world were flown at ticketed air meets --- Reims, Belmont Park,
the Gordon Bennett Cup --- and the Schneider Trophy seaplane races of the 1920s and 30s drove real
engineering progress in engines and airframes while spectators watched from the beach. If you
wanted to know the state of the art in 1925, you bought a programme. The Supermarine floatplanes
that won the Schneider led more or less directly to the Spitfire, and everybody could see them
coming.

Two things ended that. The first was the aircraft becoming a strategic weapon, which happened
between the world wars, and the second was the aircraft becoming an \emph{intelligence} weapon,
which happened after 1945. Strategic weapons need secrecy about numbers and capability.
Intelligence platforms need secrecy about existence.

The American institutional answer developed in stages. McCook Field in Dayton, Ohio, became the
Army's engineering centre in 1917; it moved to Wright Field in 1927 and became Wright-Patterson,
which is why Section~\ref{sec:wrightpattersonairforceb} keeps coming up in this book --- for decades, if the United States
acquired a piece of foreign or unexplained hardware, Dayton is where it went to be taken apart. But
Ohio has weather, terrain, population, and a fence a determined person can stand outside of. Flight
test needs something else: enormous unobstructed emergency landing area, 300-plus clear-air days a
year, no overflight of towns, and nobody within camera range.

So in 1942 the Army put its experimental flight-test operation at Muroc Dry Lake in the Mojave, on
a playa so flat and hard that a stricken aircraft could put down anywhere for miles. Muroc became
Edwards Air Force Base, and the X-plane era happened there: the X-1 breaking Mach~1 in 1947, the
D-558, the X-15 to the edge of space in the early 1960s. Notice the pattern that establishes ---
the desert dry lake as the natural habitat of the impossible aircraft. Groom Lake, when Kelly
Johnson and Richard Bissell went looking for a U-2 site in April 1955, was chosen on exactly the
same criteria as Muroc, plus one more: it already sat inside the Nevada Proving Grounds, which
meant the airspace was closed and the security apparatus was already built and paid for by the
nuclear programme. They did not have to invent secrecy. They inherited it.

And here is the consequence that matters for the rest of this book. From roughly 1955 onward there
are two histories of aviation running in parallel. There is the public one in Section~\ref{sec:historyaviation}, where
the fastest and highest thing flying is a known aircraft with a published service ceiling. And
there is the classified one, in which the actual state of the art is fifteen to twenty-five years
ahead of the published state of the art and is being flown at night over the American Southwest by
people who will deny it on the record. The U-2 was operational three years before Gary Powers made
it public by being shot down. The SR-71 was flying in 1964 and its true performance figures stayed
classified for decades. The F-117 flew combat missions in Panama before the public knew its shape.

\begin{funfact}
The gap is not a conspiracy-theory estimate; you can measure it from declassification dates. HAVE
BLUE first flew in December 1977 and was declassified in 1988 --- eleven years. Tacit Blue flew 135
sorties between 1982 and 1985 and was not revealed until 1996 --- fourteen years. The Boeing Bird of
Prey flew from 1996 and surfaced in 2002. The RQ-170 Sentinel was photographed in Afghanistan
years before it was admitted, and only because Iran captured one. So a reasonable working figure for
the classified-to-public lag in American military aviation is somewhere between ten and twenty
years. Whatever is flying out of Nevada tonight, you will read about it around 2045.
\end{funfact}

Which is why I insist that students of this subject learn aviation history properly, and not just
the famous version. A person who knows what a Skyhook balloon looks like at 90,\!000 feet at sunset,
or how a lenticular contrail forms, or what a Blackbird's plume looks like from directly beneath,
has removed most of the mystery from their own sky. A person who does not know these things is
reporting genuine observations with an empty reference library, and will reach for the only
explanation they have ever been given a picture of.

% ---------------------------------------------------------------
\section{Stealth Aircraft}\label{sec:stealthaircraft}

If you want to understand why the skies over Nevada were full of impossible objects in the
1950s, 60s, 70s and 80s, the answer is not complicated and it is not extraterrestrial. It
is that the United States was flight-testing aircraft two full generations ahead of
anything the public knew existed, in the same airspace, at night.

The sequence is worth learning because it is the single most powerful conventional
explanation in this entire field. The U-2, first flown at Groom Lake in August 1955,
cruised above 70,\!000 feet at a time when commercial airliners flew at 20,\!000 and
military aircraft rarely exceeded 40,\!000. Its silver wings caught sunlight at altitudes
where the ground was already dark. The CIA's own declassified history states plainly that
U-2 flights accounted for a substantial share of UFO reports in the late 1950s, and that
Blue Book investigators were fed deliberately false explanations --- natural phenomena,
weather inversions --- to close the files.

Then came the A-12 and SR-71 Blackbird: titanium, Mach 3.2, 85,\!000 feet, an airframe
that leaked fuel on the ground because it was designed to seal only once heated by
friction. Then HAVE BLUE and the F-117 Nighthawk, first flown in 1981 and unacknowledged
until 1988, whose faceted geometry produced a radar return the size of a ball bearing and
whose silhouette against a night sky is a flat black triangle with no visible engines.
Then Tacit Blue, an object so structurally absurd it was nicknamed the Whale.

\figwrap[l]{0.40}{The Lockheed F-117 Nighthawk. Flat, faceted, matte black, and
unacknowledged for seven years after entering service.}{https://commons.wikimedia.org/wiki/File:F-117_Nighthawk_Front.jpg}{fig:f117}
Hold that F-117 silhouette in mind, because it matters enormously for
Chapter~\ref{ch:earthfiles} and for the Belgian and Phoenix triangle waves. A large silent
flat-bottomed triangle with lights at the corners is not evidence of anything exotic. It
is a description of hardware that demonstrably existed and was demonstrably being flown
over populated areas by governments that denied its existence.

\begin{funfact}
The F-117's development was concealed by an elegant piece of misdirection: the Air Force
did not deny that something was being flown at Groom Lake. It allowed the rumour to
circulate that the secret project was a captured or reverse-engineered alien craft. A
witness who reports a spacecraft is a witness nobody in Congress will take seriously.
UFO folklore has, on at least one documented occasion, functioned as active counter
intelligence.
\end{funfact}

\begin{funfact}
Secrecy at Groom Lake was defeated, at least once, by thermodynamics. The rule was simple: when a
Soviet reconnaissance pass was due overhead, the aircraft went into the hangar. On one occasion a
U-2 had been left standing out on the ramp in the desert sun far too long before being wheeled
inside --- and the airframe had spent those hours shading the ground beneath it. Sand exposed to
full desert sun and sand kept in shadow do not return to the same temperature at anything like the
same rate, so when the overflight came, the aeroplane was safely indoors and its
\hi{outline was still lying on the apron}: a cool, sharply bounded silhouette of a wingspan and a
fuselage, invisible to the eye and perfectly legible to an infrared sensor. The Soviets
photographed the shadow of an aircraft that was not there. The story circulates in Skunk Works
lore rather than in a declassified document, so hold it loosely --- but the physics is
unimpeachable, and the moral is exact. \hi{You cannot hide an object by hiding the object.} You
have to hide everything it has touched.
\end{funfact}

Which forces an uncomfortable question on anyone honest. If the historical record shows
that the strangest objects in the sky were ours, and that the secrecy surrounding them was
total, then the null hypothesis for any given unexplained aerial object over a testing
range is: it is a black programme. Not a probe. Not a visitor. A budget line nobody will
admit to. That should be your starting position, and I hold it myself.

% ---------------------------------------------------------------
\section{Wright Patterson Air Force Base}\label{sec:wrightpattersonairforceb}

Every mythology needs a warehouse, and for American Ufology the warehouse is Hangar 18 at
Wright-Patterson Air Force Base near Dayton, Ohio.

\hi{The base is real and important, which is exactly why the story attached itself there.}
Wright Field was the technical intelligence centre of the US Air Force; the Air Technical
Intelligence Center at Wright-Patterson ran Project Sign, Project Grudge and Project Blue
Book in succession from 1948 to 1969. If any physical material from any American UFO case
was ever analysed by the government, Wright-Patterson is the correct place to look. It was
the institutional destination for captured and unusual foreign hardware --- German,
Japanese, later Soviet. That is documented, mundane, and entirely sufficient to explain
why debris of any kind would be shipped there.

The claim, which is not documented, is that a hangar at the base holds recovered
extraterrestrial craft and bodies. The story attaches most persistently to Roswell
material and to the alleged testimony of Senator Barry Goldwater, who stated in
correspondence that he had asked General Curtis LeMay for permission to enter a room at
Wright-Patterson known as the ``Blue Room'' and was refused in extremely forceful terms.
Goldwater's account is genuine --- the letters exist and he repeated the story publicly ---
but note carefully what it establishes: that a US senator was denied access to a secure
room. That is not evidence about the room's contents. It is evidence about compartmented
clearance, which functions exactly that way by design.

As for ``Hangar 18'' specifically: the base has no building with that designation, and the
number appears to originate in 1970s paperback literature and a 1980 film of that title.
This is the Section~\ref{sec:confirmationbias} laundering problem in its purest form. Trace the term and the trail
ends at a novel.

What \emph{is} solidly documented at Wright-Patterson is duller and more damning. The Blue
Book files were kept there. Their organisation was poor, their staffing was minimal --- for
long stretches a single officer and a secretary --- and case dispositions were frequently
recorded with no investigation whatsoever. If you want to be angry about
Wright-Patterson, be angry about that. The scandal is not a hidden hangar. The scandal is
an open filing cabinet full of unexamined reports.

% ---------------------------------------------------------------
\section{JANET Airlines}\label{sec:janet}

Here is my favourite fact in this entire chapter, and it is not a rumour, a leak, or a
deathbed confession. \hi{The most secret facility in the United States has a scheduled
commuter airline, and you can go and watch it operate.}

\figwrap[r]{0.40}{A \term{JANET} Boeing 737-600 in the fleet's white livery with its single
red cheatline. No airline name, no operator marking.}{https://commons.wikimedia.org/wiki/File:Janet_Boeing_737-66N.jpg}{fig:janet737}
On the western side of Harry Reid International Airport in Las Vegas --- McCarran, when I
first went looking --- there is a fenced compound with its own car park, its own gate, and
its own apron. It is officially the Gold Coast Terminal. Everyone calls it the Janet
terminal. Several times a morning, a white Boeing 737 with a single red stripe down the
fuselage and no airline name anywhere on it taxis out, takes off, and flies north into
restricted airspace. Several times an afternoon, one comes back. The passengers are
cleared employees of the Nevada National Security Site, the Tonopah Test Range and Groom
Lake, and they commute to work the way accountants in Connecticut commute to Manhattan.

The operation is old. It began in 1972 with a Douglas DC-6B run by EG\&G, the defence
contractor that has managed Nevada test-site support work for most of the nuclear age, and
it has changed hands with the support contract ever since; the fleet is currently operated
under contract by Amentum. The airliners today are six second-hand Boeing 737-600s,
acquired in 2008, which had previously flown for Air China --- an aircraft type retired by
almost every commercial operator on the planet, which is part of why the fleet is so easy
to spot. Alongside them sits a handful of smaller Beechcraft turboprops used for shorter
hops within the range complex.

The call signs are where it becomes properly entertaining. The aircraft file under the
radio call sign ``Janet'' followed by a three-digit number, and their registrations all
begin with N and the letters WW. Air-traffic recordings show the flights switching call
sign as they enter the restricted airspace: ``Janet'' becomes ``Groom'', and the numbers
shift by a fixed offset. The destination itself is not called Area 51 on the paperwork.
In the flight-crew documentation it is Station 3. None of this is secret in any operational
sense; it is on public frequencies, and hobbyists have logged it for decades.

\begin{funfact}
Nobody has ever confirmed what ``Janet'' stands for. The tidiest theory, and the one you
will see repeated as fact, is ``Joint Air Network for Employee Transportation''. That
phrase looks exactly like what it probably is: a backformation invented to fit an existing
word. Competing folklore holds that a Groom Lake controller simply named the flights after
his wife. I have no idea which is true, and I want to flag that ignorance rather than
smooth it over, because this is a case where the mundane detail is genuinely unrecorded
while the exotic details around it are documented to the tail number.
\end{funfact}

So what does JANET actually prove? Two things, and they cut in opposite directions, which
is why I have given it a section of its own.

First, it demolishes a certain kind of maximalist secrecy claim. A facility that requires
somewhere between one and two thousand seats of daily airline capacity is not a small
crew of technicians guarding nine flying saucers in a mountain. It is a large, industrial,
unionised workplace with shift patterns, a car park problem and a lost-luggage procedure.
The scale argues strongly for exactly what the documented history says the base does:
flight-testing expensive aircraft, which requires engineers, machinists, range safety
staff and administrators in industrial quantities. Secrets kept by thousands of commuters
leak, and what has leaked over fifty years is aircraft.

Second, and less comfortably, it shows how thin the line between ``classified'' and
``obvious'' really is. The government spent decades refusing to admit Groom Lake existed
while running a visible airline to it from a civilian international airport. Anyone with a
long lens and a parking ticket could document the traffic. That is not competent secrecy;
it is compartmented secrecy, which protects \emph{what} is done at the base rather than
\emph{that} the base is there. Understanding that distinction is one of the more useful
tools in this whole subject, because a great deal of Ufology consists of discovering the
visible half of a compartmented programme and concluding that the invisible half must
therefore be extraordinary.

One further note, because it matters and it is usually left out of the fun version of this
story. On 12 January 2004, a Beech 1900C on a Janet-operated flight crashed shortly after
taking off from the Tonopah Test Range. All five people aboard were killed. The commute is
a real job with real risk, performed by people who cannot tell their families what they did
that day. When you photograph the 737s from the garage roof --- and you should, it is one of
the strangest legal activities in American aviation --- it is worth remembering that the
passengers are not props in a mystery.

\begin{srcnotes}
\item ``Janet (airline)'', Wikipedia --- fleet history, registrations, call-sign practice and
terminal details: \srclink{https://en.wikipedia.org/wiki/Janet_(airline)}.
\item Federal Aviation Administration registry entries for the WW-series N-numbers
(publicly searchable).
\item National Transportation Safety Board report on the 12 January 2004 Beech 1900C
accident at Tonopah Test Range.
\item Jacobsen, A.\ (2011). \emph{Area 51: An Uncensored History}. New York: Little, Brown
--- on EG\&G's role as the site's long-term support contractor.
\end{srcnotes}

% ---------------------------------------------------------------
\Needspace*{0.40\textheight}
\section{S4 and Bob Lazar}\label{sec:s4lazar}

\figfull[3.4in]{Papoose Lake, Nevada, the dry lakebed south of Groom Lake where Lazar placed the S-4 hangars. No photograph of the ``sport model'' disc exists, and Lazar's own sketch of it is not free to reproduce, so the site is shown instead: the terrain is documented, the facility in the mountainside is not.}{https://commons.wikimedia.org/wiki/File:Papoose_Lake_99-7621_(5881444036).jpg}{fig:lazardisc}
In May 1989, a man appeared in silhouette on Las Vegas television station KLAS under the
pseudonym ``Dennis'' and told reporter George Knapp that he had been employed as a
physicist at a facility called S-4, at Papoose Lake south of Groom Lake, reverse
engineering the propulsion system of an extraterrestrial craft. He went public under his
own name later that year. His name is Robert Lazar, and the modern Area 51 mythology is
very largely his creation.

\textbf{\textcolor{ufogreendark}{His claims are specific, which is to his credit}} --- specificity is falsifiable. He described nine craft in hangars built into the base of a mountain with sloped doors, at the site shown in Figure~\ref{fig:lazardisc}. He
described a reactor running on a stable superheavy element he called Element 115, which
produced antimatter on proton bombardment and generated a gravitational field that could
be amplified and directed to warp spacetime around the craft. He described being briefed
with documents concerning long-term extraterrestrial involvement with Earth. He described
seeing a craft flight-tested.

Now let us do the work, because this case is an excellent exercise.

What holds up: Lazar's knowledge of the Nevada terrain and of the general geography of
test operations was good. He led Knapp and others to a location from which they observed
unusual aerial activity on Wednesday evenings, which he had predicted in advance. And
element 115 --- which did not exist in 1989 --- was subsequently synthesised in 2003 and
named moscovium in 2016. Lazar's supporters treat this as a bullseye.

What does not hold up is more substantial. Moscovium is not a stable superheavy element;
its longest-lived known isotope has a half-life measured in tenths of a second, and it
cannot be a fuel for anything. Predicting that element 115 would eventually be made was
not prescience --- the periodic table has numbered blanks, and physicists had been
discussing an ``island of stability'' in the popular press for years. More seriously, his
claimed education has never been substantiated. He states he holds degrees from MIT and
Caltech; neither institution has any record of him, no classmate or instructor has ever
come forward, and the employment history he has offered is documented only in fragments
that do not require any physics credential. He has a 1990 conviction connected to a
prostitution business, which is not relevant to physics but is relevant to the claim that
he was a cleared contractor.

\begin{funfact}
The most durable part of Lazar's account is the one nobody argues about: he put the phrase
``Area 51'' into general public circulation. Before the 1989 KLAS broadcasts the name was
known almost exclusively to aviation specialists. Whatever else Bob Lazar did or did not
do, he named the most famous secret in the world.
\end{funfact}

My own assessment, offered as an assessment and not a finding: the credentials do not
concern me. Jack Parsons had no credentials before he built the Jet Propulsion Laboratory.
Einstein was not educated as a mathematician before he wrote the equations that bear his name.
Credentials certify existing institutions; they do not certify curiosity, and they never have.
As for how Lazar ended up at S-4, his account in the Project Gravitaur film describes a
fairly routine ascension through the ranks --- first at EG\&G, then into contracted
operations with other government departments --- and there is nothing abnormal about it.
Workers float between contractors inside the larger framework of that machine all the time.
The how is not the mystery.

What fascinates me is the physics. Lazar described a cylindrical black barrel, hollow,
made of some kind of metal, that when rotated ninety degrees froze time in the room: a
candle flame stopped mid-air. Two things stand out. First: why ninety degrees? That is a
vector dot product --- a perfect perpendicular --- and it proves that mathematics went into
whatever the device was designed to do. Second: the barrel froze time for the flame but
not for the observers. Perhaps human beings are not as bound to time as we initially
assumed.

Moscovium remains the checkable objection. Its longest-lived isotope lasts around two
tenths of a second and cannot be a stable fuel, and that part of the account does not
survive the periodic table. I say this while noting, as
Section~\ref{sec:doubtdrivenskepticism} requires, that ``he cannot document his degrees'' does
not prove he was never at Papoose Lake, and that a government which lied about the
existence of Groom Lake for fifty years is not a reliable witness against him. The case is
not closed by official denial. It is not closed by moscovium either. The moscovium
objection closes one claim. It does not close the room where the barrel sat.
% ---------------------------------------------------------------
\section{Cultural Impact of Area 51}\label{sec:culturalimpactarea51}

Area 51 has become the standard-issue backdrop for secrecy itself, and the trajectory is
worth tracing because it shows how a place becomes a symbol and then stops being a place.

The 1989 Lazar broadcasts gave it a name. \emph{Independence Day} in 1996 gave it a
location in the popular imagination, complete with an underground hangar --- a set, not a
photograph, though a startling number of people cannot now tell you which. The Nevada
state legislature renamed State Route 375 the ``Extraterrestrial Highway'' in 1996 as a
tourism measure, and the Little A'Le'Inn at Rachel became a genuine economic
beneficiary of a rumour. The base now appears in dozens of video games as a level.

Then, in June 2019, a college student created a Facebook event titled ``Storm Area 51,
They Can't Stop All of Us'' as a joke. Two million people clicked attending. The Air Force
issued a formal statement warning that the range was a live training area and that
attempts to enter would be met with force. Two small festivals were hastily organised in
Rachel and Hiko; roughly 150 people turned up at the gates, six were arrested, and a
remote Nevada county came genuinely close to a humanitarian problem because a meme
outgrew its punchline.

Around the base a whole pilgrimage geography grew up, and I want to name its stations because
they are the physical residue of the cultural story. Nevada State Route 375 carries roughly two
hundred vehicles a day, which makes it one of the least-used numbered highways in the country and
an improbable candidate for a tourist attraction. The Little A'Le'Inn opened at Rachel in 1989 ---
the same year as the Lazar broadcasts, which is not a coincidence --- and later turned up in the
film \emph{Paul}. Twelve miles from the Groom Lake road turnoff stood the Black Mailbox, the
property marker of a rancher named Steve Medlin and the only landmark of any kind within fifty
miles. It became a gathering point for night-time skywatchers, and Medlin eventually replaced it
with a white one in the hope that people would stop finding it; they did not, and he ultimately
listed the property for sale. Before the 1984 and 1995 land seizures, an amateur group known as the
Skywatchers Association camped on Whitesides and Freedom Ridge, watched closely by base security
from the other side of the line. Since those vantage points were taken, the only remaining spot from
which the base is visible to the naked eye is Tikaboo Peak, twenty-six miles out --- roughly the
distance from Lethbridge to Fort Macleod, if that helps you picture it --- and reaching it is a
serious hike with a telescope on your back. In 2012 the National Atomic Testing Museum in Las Vegas,
a Smithsonian affiliate, opened an Area 51 exhibit, which is a small institutional milestone worth
noting: the artefacts of the rumour are now curated inside a federally-affiliated museum, a year
before the CIA admitted the base existed.

I find the whole arc instructive rather than annoying. Area 51 in 2026 is not primarily a
Ufological subject. It is a case study in how official silence becomes a canvas. The
government's refusal to acknowledge the base did not protect the secret --- the aircraft
were all eventually revealed, on the government's own schedule, with fanfare. What the
silence produced was a vacuum, and vacuums get filled by whoever is willing to speak. In
this instance that was a man with unverifiable degrees and a television reporter with an
hour to fill, and their version has outlasted every declassification since.

There is a lesson in there for anyone in a position of institutional authority who
imagines that saying nothing is the same as saying nothing happened. It is not. \hi{Saying
nothing is an invitation, and somebody always accepts.}

\begin{funfact}[Fun Fact: the fax from the wildcard line]
One of the strangest fragments in the Area 51 folklore was not a document or a photograph. It was a
fax. On \emph{Coast to Coast AM} --- Art Bell's programme, discussed in Section~\ref{sec:artbell} ---
Bell read out a message that had just come in: \emph{``Art! Answer the wildcard line. My boyfriend is
in a small plane north of Las Vegas. Ready to fly into Area 51. He's been trying to call you by
cellphone.''} Bell took the call. The man was audibly airborne, engine noise behind his voice, and he
narrated his approach on air until the transmission cut off mid-sentence. In the version that has
circulated ever since, he was shot down.

I include this because it is a perfect specimen of how the mythology is manufactured, and because
honesty requires me to say what is actually established. \hi{The only detail I can verify from the
recording itself is the engine noise} --- that is, that the caller was in an aircraft of some kind. I
have not been able to establish his name, the date of the broadcast, the aircraft, or any
corroborating record of an interception, a crash, or a fatality. A cut-off cellphone call from a light
aircraft over the Nevada desert in the 1990s has a mundane explanation available at all times. Note
also that this is a separate episode from the far more famous ``Area 51 employee'' call of 11
September 1997, which is often confused with it. The story survives not because it was ever
documented, but because the recording ends at exactly the right moment --- and, as the rest of this
chapter argues, a silence is the most persuasive thing a mystery can offer.
\end{funfact}

\begin{srcnotes}
\item \emph{Coast to Coast AM} broadcast audio, host Art Bell, reading the ``wildcard line'' fax and
taking the airborne caller. Circulating copies of the clip give no broadcast date, and the version
formerly hosted on YouTube is no longer available; the quotation above follows the fax text as read on
air.
\item For the distinct and much better documented ``Area 51 employee'' call of 11 September 1997, and
for Bell's broadcasting career generally, see Chapter~\ref{ch:media}, Section~\ref{sec:artbell}.
\end{srcnotes}
